#!/usr/bin/env python3
"""
convert_papers.py
=================
Converts a plain-text MathSciNet bibliography file into a tab-separated
file suitable for David Eisenbud's papers list.

INPUT FORMAT (MathSciNet plain text):
    MR4968625  Author1, First ; Author2, First . Title of paper.
     Journal Name  VOL  (YEAR),  no. N, PAGE--PAGE. YEAR-NNN

OUTPUT FORMAT (9 tab-separated columns):
    INDEX  YEAR  TITLE  AUTHORS  JOURNAL  VOL:ISSUE  PAGES  YEAR-CODE  ?

USAGE:
    python3 convert_papers.py input.txt output.txt [--start N]

    --start N   : starting index number (default: 1)

NOTES:
- Entries in the input file are processed in REVERSE order (last entry
  gets the lowest index number, matching the existing list convention).
- The MR token is stripped from each entry.
- Book chapters: journal column gets the book/series title; vol/issue blank.
- The input file should have Unix line endings and a final newline.
  If needed, preprocess with:
      sed -i '' 's/--/-/g' input.txt   (replace double hyphens)
      printf "\\n" >> input.txt         (ensure final newline)
- The downstream awk script requires /regex/ syntax (not "string") in
  sub() calls to avoid '.' matching any character. Specifically:
      sub(/\\.$/, "", author)       -- not sub("\\.$", ...)
      sub(/\\.\\.$/,  ".", title)   -- not sub("\\.\\.$", ...)
"""

import re
import argparse
import sys

# -----------------------------------------------------------------------
# Keywords that identify book chapters / conference proceedings.
# If any of these appear in the publication info, the entry is treated
# as a book chapter: col5 = book/series title, col6 = blank.
# Extend this list as needed.
# -----------------------------------------------------------------------
BOOK_KEYWORDS = [
    'Contemp. Math.',
    'London Math. Soc.',
    'Abel Symp.',
    'Clay Math.',
    'Math. Sci. Res. Inst.',
    'Hindustan Book',
    'Cambridge Univ. Press',
    'Springer, Berlin',
    'Amer. Math. Soc.',
    'Proceedings of the International Congress',
]

# -----------------------------------------------------------------------
# Manual fixes: applied AFTER automatic parsing.
# Keys are 1-based positions in the REVERSED entry list (i.e. the output
# row number relative to --start).
# Use this for entries where the parser can't cleanly separate journal
# name from volume (usually because source has only single spaces).
# Format: position -> (journal, vol_issue, pages)
#
# When reusing this script for a new batch, clear this dict and add new
# fixes as needed after a first-pass review.
# -----------------------------------------------------------------------
MANUAL_FIXES = {
    # Examples from the 2007-2025 batch (positions relative to start=134):
    # 134: ('Bull. Amer. Math. Soc. (N.S.)', '44:3', '331-359'),
    # Add new fixes here as needed.
}


def parse_authors(s):
    """Split 'Last, First ; Last, First ; ...' into 'Last, First, Last, First, ...'"""
    parts = [p.strip().rstrip('.') for p in re.split(r'\s*;\s*', s)]
    return ', '.join([p for p in parts if p])


def parse_entry(raw):
    """
    Parse one raw MathSciNet entry (MR token already present) into a dict
    with keys: year, title, authors, journal, vol_issue, pages, year_code.
    """
    # --- Strip MR token ---
    raw = re.sub(r'^MR\d+\s*', '', raw)

    # --- Extract year-code (e.g. 2025-001) from end of entry ---
    m = re.search(r'(\d{4}-\d{3})\s*$', raw)
    year_code = m.group(1) if m else '?'
    raw = raw[:m.start()].strip() if m else raw
    year = year_code[:4]

    # --- Split into lines, clean up ---
    lines = [l.strip() for l in raw.split('\n') if l.strip()]
    # Remove language annotations like "(French)"
    lines = [re.sub(r'^\([A-Z][a-z]+\)\s*', '', l) for l in lines]

    # ---------------------------------------------------------------
    # Parse authors and title from line 0.
    # MathSciNet format: "Author1 ; Author2 . Title text."
    # The separator is ' . ' (space-dot-space).
    # Special case: some entries have "Nitsche, Max J.  Title" where
    # the author list ends with an initial and double-space precedes title.
    # ---------------------------------------------------------------
    line0 = lines[0]
    sep = re.search(r'\s\.\s', line0)
    if sep:
        authors_raw = line0[:sep.start()]
        title = line0[sep.end():].strip().rstrip('.')
    else:
        # Fallback: author block ends with "X." then double space
        m2 = re.match(r'^((?:[^.]+\.\s*)+?)\s{2,}(.+)$', line0)
        if m2:
            authors_raw = m2.group(1).strip()
            title = m2.group(2).strip().rstrip('.')
        else:
            authors_raw = line0.rstrip('.')
            title = ''

    # --- Check if title continues on line 2 (lowercase continuation) ---
    pub_start = 1
    while pub_start < len(lines):
        l = lines[pub_start]
        if l and l[0].islower():
            # Lowercase start = continuation of title
            title += ' ' + l.rstrip('.')
            pub_start += 1
        elif re.match(r'^Special volume', l):
            # Dedication note (e.g. "Special volume in honor of X") — skip
            pub_start += 1
            break
        else:
            break

    # --- Clean up title ---
    authors = parse_authors(authors_raw)
    title = re.sub(r'\s+', ' ', title).strip()
    title = title.replace('\\ell', 'ell')          # LaTeX \ell -> ell
    title = re.sub(r'\[\s*MR\d+\s*\]', '[MR]', title)  # shorten MR refs

    # --- Publication info: join remaining lines ---
    pub_text = ' '.join(lines[pub_start:])
    pub_text = re.sub(r'\s+', ' ', pub_text).strip()
    # Remove any dedication note at start of pub_text
    pub_text = re.sub(r'^Special volume[^.]+\.\s*', '', pub_text)

    # ---------------------------------------------------------------
    # Parse publication info into journal, vol_issue, pages.
    # Three cases:
    #   1. Book chapter (contains BOOK_KEYWORDS)
    #   2. Épijournal-style with [year range]
    #   3. Standard journal article
    #      a. Gaz. Math. style: "Journal  No. N,  (YEAR), pages"
    #      b. Normal: "Journal  VOL  (YEAR),  [no. N,]  pages"
    # ---------------------------------------------------------------
    is_book = any(kw in pub_text for kw in BOOK_KEYWORDS)
    is_epi  = 'pijournal' in pub_text or ('[' in pub_text and '--' not in pub_text
                                          and re.search(r'\[\d{4}', pub_text))

    journal = vol_issue = pages = ''

    if is_book:
        pm = re.search(r'(\d+-\d+)', pub_text)
        pages = pm.group(1) if pm else ''
        if pages:
            book_part = pub_text[:pub_text.index(pages)].strip().rstrip(',').strip()
        else:
            book_part = pub_text.split(',')[0].strip()
        journal = book_part

    elif is_epi:
        m = re.match(r'^(.*?)\s+\[.*?\],\s*(.*)', pub_text)
        if m:
            journal = m.group(1).strip()
            rest = m.group(2)
            art = re.search(r'Art\.\s*(\d+)', rest)
            pp  = re.search(r'(\d+)\s*pp', rest)
            if art and pp:
                pages = f"Art. {art.group(1)}, {pp.group(1)} pp"
            elif pp:
                pages = pp.group(1) + ' pp'
        else:
            journal = pub_text

    else:
        # Gaz. Math. style: no volume number, just "No. N"
        gaz = re.match(r'^(.*?)\s+No\.\s*(\d+),\s*\(\d{4}\),\s*(.*)', pub_text)
        if gaz:
            journal   = gaz.group(1).strip()
            vol_issue = 'No. ' + gaz.group(2)
            rest = gaz.group(3).rstrip('.')
            pm = re.search(r'(\d+-\d+)', rest)
            pages = pm.group(1) if pm else rest

        else:
            # Standard journal: try double-space split first, then single-space
            m = re.match(r'^(.*?)\s{2,}(\d+)\s+\(\d{4}\),\s*(?:no\.\s*(\S+),\s*)?(.*)', pub_text)
            if not m:
                m = re.match(r'^(.*?)\s(\d+)\s+\(\d{4}\),\s*(?:no\.\s*(\S+),\s*)?(.*)', pub_text)
            if m:
                journal   = m.group(1).strip()
                vol       = m.group(2)
                issue     = m.group(3)
                rest      = m.group(4).strip().rstrip('.')
                vol_issue = f"{vol}:{issue}" if issue else vol
                pm = re.search(r'(\d+-\d+)', rest)
                if pm:
                    pages = pm.group(1)
                else:
                    pp = re.search(r'(\d+)\s*pp', rest)
                    pages = pp.group(1) + ' pp' if pp else rest
            else:
                journal = pub_text

    return dict(year=year, title=title, authors=authors,
                journal=journal, vol_issue=vol_issue, pages=pages,
                year_code=year_code)


def main():
    parser = argparse.ArgumentParser(description=__doc__,
                                     formatter_class=argparse.RawDescriptionHelpFormatter)
    parser.add_argument('input',  help='Input plain-text file from MathSciNet')
    parser.add_argument('output', help='Output tab-separated file')
    parser.add_argument('--start', type=int, default=1,
                        help='Starting index number (default: 1)')
    args = parser.parse_args()

    with open(args.input, 'r', encoding='utf-8') as f:
        text = f.read()

    # Split on blank line before each MR token
    entries = re.split(r'\n(?=MR\d+)', text.strip())
    entries = [e.strip() for e in entries if e.strip()]

    # Reverse: last entry in file gets the lowest index number
    entries = list(reversed(entries))

    results = [parse_entry(e) for e in entries]

    # Apply manual fixes
    for i, r in enumerate(results):
        pos = args.start + i
        if pos in MANUAL_FIXES:
            r['journal'], r['vol_issue'], r['pages'] = MANUAL_FIXES[pos]

    # Write output
    output_lines = []
    for i, r in enumerate(results):
        num  = args.start + i
        cols = [str(num), r['year'], r['title'], r['authors'],
                r['journal'], r['vol_issue'], r['pages'], r['year_code'], '?']
        output_lines.append('\t'.join(cols))

    with open(args.output, 'w', encoding='utf-8') as f:
        f.write('\n'.join(output_lines) + '\n')

    print(f"Wrote {len(output_lines)} entries to {args.output}.")

    # Warn about any entries that may need manual review
    issues = []
    for i, r in enumerate(results):
        num = args.start + i
        if not r['title']:
            issues.append(f"  Entry {num}: empty title")
        if re.search(r'\(\d{4}\)', r['journal']):
            issues.append(f"  Entry {num}: year found in journal field — may need manual fix")
    if issues:
        print("\nEntries that may need review:")
        for issue in issues:
            print(issue)
        print("\nAdd corrections to the MANUAL_FIXES dict at the top of the script.")


if __name__ == '__main__':
    main()
